\scanpage{6}
\mspage{3}
\[
(14)\qquad1\mathbin{\unsure{-|\to}}\pi_1(X')\mathbin{\unsure{-|\to}}\mathcal E_{X'}\longrightarrow\mathcal G_{X'}\longrightarrow1.
\]
et on a encore un germe de relèvement
(\unsure{éventuellement} mod comp. par autom int.
provenant de $\pi_1(X')$)
\[
(15)\qquad\diagram{\begin{tikzcd}[ampersand replacement=\&]
\mathcal E_{X'}\arrow[r,dash]\&\mathcal G_{X'}\arrow[r]\&\Gamma\arrow[d,phantom,"\cup" description]\\
\&\&\Gamma'\arrow[ull,"\varphi_\xi"']
\end{tikzcd}}
\]
(si $\xi\in X'(\overline{\mathbf Q})$).

Notons qu'on a ici un hom
\[
(16)\qquad\mathcal G_{X'}\longrightarrow\mathfrak S_I
\]
La théorie de Kummer nous fournit,
pour $\forall\,i$, un hom extérieur
\[
(17)\qquad\mathbf T(\overline{\mathbf Q})\xrightarrow{\varphi_i}\pi_1(X')
\]
\[
=\varprojlim\mu_n(\overline{\mathbf Q})
\]
\margin{$\varphi_i\in\operatorname{Homext}(\mathbf T(\overline{\mathbf Q}),\pi_1(X'))$}
\margin{\deleted{\unsure{il faut supposer} $\chi=1$ \unsure{si l'on veut}
\unsure{interpréter ça par} \illegible{6} \unsure{purement arithmétique} \illegible{3}}}
Notons aussi un hom canonique
\[
(18)\qquad(\mathcal G_{X'}\mathbin{-})\Gamma\xrightarrow{\chi}\widehat{\mathbf Z}^{*}\xrightarrow{\sim}\operatorname{Aut}(\mathbf T(\overline{\mathbf Q}))
\]
faisant opérer $\mathcal G_{X'}$ via $\Gamma$ sur $\mathbf T(\overline{\mathbf Q})$.
Ceci dit, l'opération extérieure de $\mathcal G_{X'}$
sur $\pi_1(X')$ est compatible avec la famille
\mspage{5}
